\documentclass[a4paper,12pt]{article}
\usepackage[english,russian]{babel}
\usepackage{fontspec}
\defaultfontfeatures{Ligatures={TeX},Renderer=Basic}
\setmainfont[Ligatures={TeX,Historic}]{Calibri Light}
\setsansfont{Calibri Light}
\setmonofont{Consolas}

\usepackage{indentfirst}
\frenchspacing

% Для математики
\usepackage{amssymb,amsmath}
\parindent=24pt
\parskip=0pt
\tolerance=2000

% Для настройки размера страницы
\usepackage{geometry}
\geometry{
	a4paper,
	total={170mm,257mm},
	left=20mm,
	top=20mm,
}

% Для вставки графики
\usepackage{graphicx}

% Для таблиц
\usepackage{multirow}
\usepackage{tabu}

% Создаем команду, чтобы переносить текст на новую строку внутри таблицы
\newcommand{\tcell}[2][c]{\begin{tabu}[#1]{@{}c@{}}#2\end{tabu}}

% Пакет для списков
\usepackage[ampersand]{easylist}

% Для контура вокруг текста
\usepackage[outline]{contour}

% Пакеты от tikz
\usepackage{tikz}
\usepackage{graphics}
\usepackage{pgfplots}
\usepackage{pgfplotstable}
\usepackage{xcolor}
\usetikzlibrary{calc}
\usetikzlibrary{through}
\usetikzlibrary{intersections}
\usetikzlibrary{patterns}
\usetikzlibrary{scopes}
\usetikzlibrary{decorations.pathreplacing}
\usetikzlibrary{arrows.meta}

% Для цветных таблиц
\usepackage{tcolorbox}
\tcbuselibrary{skins}
\tcbuselibrary{breakable}
\tcbuselibrary{minted}
\usemintedstyle{vs}

% Для подсветки кода
\usepackage{minted}

% Задаем цвет номеров строк подсветки кода
\renewcommand{\theFancyVerbLine}{\sffamily\textcolor[rgb]{1, 1, 1}{\scriptsize\arabic{FancyVerbLine}}}

% Определяем новую команду для красивой вставки кода в рамочке и с прочими прелестями
\newcommand{\mycodeinput}[3]{
\begin{tcolorbox}[
	colback=black!5!white,
	colframe=black!30!white,
	boxrule=0.5pt, 
	listing only,
	left=-0.5mm,
	leftrule=6mm,
	arc=2mm, outer arc=2mm,
	top=2pt,
	bottom=0pt,
	enhanced,
	breakable,
	title={#3},
	coltitle=black, 
	fonttitle=\bfseries\ttfamily\large,
	attach boxed title to top center={yshift=-3mm,yshifttext=-1mm},
	boxed title style={
		enhanced,
		nobeforeafter,
		tcbox raise base,
		boxrule=0.4pt,
		top=0.5mm,
		bottom=0.5mm,
		right=0mm,
		left=4mm,
		arc=1pt,
		boxsep=3pt,
		before upper={\vphantom{dlg}},
		colframe=black!30!white,
		colback=black!10!white,
		overlay={
			\begin{tcbclipinterior}
			\fill[black!30!white]
				(frame.south west)
					rectangle node[text=white,font=\sffamily\bfseries\footnotesize,rotate=90] {FILE} 
				([xshift=4mm]frame.north west);
			\end{tcbclipinterior}
		}
	},
]
\inputminted[
	breaklines,
	breakanywhere=true,
	autogobble,
	linenos,
	numbersep=1mm,
	mathescape, 
	fontsize=\footnotesize, 
	tabsize=4
]{#1}{#2}
\end{tcolorbox}
}

% Определяем новую команду для создания титульного листа
\newcommand{\mytitlepage}[9]{
\begin{center}
\hfill \break
\Large{Министерство образования и науки Российской Федерации}\\
\hfill \break
\large{Федеральное государственное бюджетное образовательное учреждение высшего образования}\\ 
\normalsize{\textbf{«НОВОСИБИРСКИЙ ГОСУДАРСТВЕННЫЙ ТЕХНИЧЕСКИЙ УНИВЕРСИТЕТ»}}\\
\hfill \break
\includegraphics{nstu_logo.eps} \\
\hfill \break
\large{Кафедра #1}\\
\hfill \break
\large{Лабораторная работа №#2\\по дисциплине <<#3>>}\\
\hfill \break
\hfill \break
\Large{\textbf{#4}}\\
\hfill \break
\hfill \break 
\normalsize{\begin{tabular}{cllp{1.5cm}p{1.5cm}}
\multirow{5}{*}[0.75cm]{\includegraphics[scale=0.5]{fami_logo.eps}}
& \textbf{Факультет:} & ПМИ & & \\[1.25ex]
& \textbf{Группа:} & #5 & & \\[1.25ex]
& \textbf{Студент:} & #6 & & \\[1.25ex]
& \textbf{Вариант:} & #7 & & \\[1.25ex]
& \textbf{Преподаватель:} & #8 & & \\[1.25ex]
\end{tabular}} \\
\hfill \break
\hfill \break
\hfill \break
\hfill \break
\hfill \break
\large{Новосибирск\\#9}
\end{center}
\thispagestyle{empty}
\newpage 
\setcounter{page}{1}
}

\begin{document}

\mytitlepage{прогресса}{1}{Создание отчетов}{Создание отчета в LaTeX}{ПМ-63}{Шепрут И.И.}{1}{Шепрут И.И.}{2018}

\section{Цель работы}

Показать пример использования LaTeX для написания отчетов.

\section{Анализ задачи}

Для этого необходимо собрать минимальное необходимое количество примеров, а так же оставить ссылки на документацию.

\section{Код программы}

Для демонстрации этого написан текущий файл с таким кодом:

\mycodeinput{latex}{my_title.tex}{title.tex}
\mycodeinput{latex}{file.tex}{file.tex}

\section{Таблица}

\tabulinesep=0.3mm
\noindent\scriptsize{\texttt{\begin{tabu}{|X[-1,c]||X[1,c]|X[1,c]|X[1,c]|X[1,c]|X[1,c]|}
\hline
	\# & Название & Документация & Пример & PDF & Присутствует в этом файле \\
\hline
\hline
	1 & Code & + & + & + & + \\
\hline
	2 & Graph & + & + & + & + \\
\hline
	3 & Graph by data & + & + & + & $\pm$ \\
\hline
	4 & List & + & + & + & + \\
\hline
	5 & Table & + & + & + & + \\
\hline
\end{tabu}}}

\section{График}

\noindent\begin{tikzpicture}
	\begin{axis}[
		xlabel=Пример,
		ylabel=Присутствует в этом файле?,
		width=\textwidth, 
		height=5cm]
	\addplot[color=red, mark=x] coordinates {
		(1,1)
		(2,1)
		(3,0.5)
		(4,1)
		(5,1)
	};
	\end{axis}
\end{tikzpicture}

\section{Пример формулы}

$$ f(x) = \sum_{k=0}^n \frac{f^{(k)}(x_0)}{k!}(x-x_0)^k + \underbrace{\frac{1}{n!}\int\limits_{x_0}^x f^{(n+1)}(t)(x-t)^n dt}_{\text{интегральный остаточный член}} $$

Код \mycodeinline{latex}{$\frac{1}{a}$} сделает дробь $\frac{1}{a}$.

\section{Выводы}

\noindent\normalsize{\begin{easylist}
\ListProperties(Hang1=true)
& Было показано множество полезных примеров.
& Документация расположена в папках примеров.
& В ворде некоторые вещи делаются в миллион раз проще и в пару кликов.
& Зато в латехе ничего не съедет, и можно автоматически генерировать таблицы.
\end{easylist}}

\end{document}